%%%PAGE 66 rot=0 legibility=good summary=匀变速直线运动章节提纲(思维导图式)：描述运动的物理量、规律、基本公式与推论的应用、自由落体竖直上抛、多过程、异表同质、图像、追及相遇；右侧附初速度为零的匀变速比例式
%%%BLOCK id=066-01 ch=01 sec=章节提纲 kind=summary alt=none cont=none y=0.07-0.97
% 版面：左侧“匀变速直线运动”后接大括号，括入前四项；其下另有一个大括号，括入“匀变速直线运动的规律”“自由落体和竖直上抛规律”；再往下各行无括号。
% “匀变速直线运动推论的应用”“自由落体和竖直上抛运动”两行右侧的比例式批注见 066-02。
\textbf{匀变速直线运动}
\begin{itemize}
  \item 质点、参考系
  \item 位移、速度
  \item 路程、速率
  \item 加速度
\end{itemize}
\begin{itemize}
  \item 匀变速直线运动的规律
  \item 自由落体和竖直上抛规律
\end{itemize}
\begin{itemize}
  \item 匀变速直线运动基本公式的应用
  \item 匀变速直线运动推论的应用
  \item 自由落体和竖直上抛运动
  \item 匀变速直线运动多过程问题\quad 找衔接速度。
  \item 异表同\unclear{质}问题
  \begin{itemize}
    \item 水平刹车与粗糙斜面上滑
    \item 竖直上抛与光滑斜面（足够长）上滑。
    \item 逆向思维法
  \end{itemize}
  \item 运动图像的理解应用
  \item 追及相遇问题
\end{itemize}
%%%END
%%%BLOCK id=066-02 ch=01 sec=匀变速直线运动推论·初速度为零的比例式 kind=knowledge alt=none cont=none y=0.50-0.66
% 原稿位于“匀变速直线运动推论的应用”“自由落体和竖直上抛运动”两行的右侧
$1T$末\ $2T$末\ $3T$末，…$nT$末瞬时速度之比
\[ v_1:v_2:v_3:\cdots:v_n=1:2:\cdots:n \]
第$1$个$T$内，$2$个$T$内，$3$个$T$内…$n$个$T$内位移之比
\[ x_1:x_2:x_3=\cdots:x_n=1:3:5:\cdots:\text{\struck{\unclear{…}}}(2n-1) \] % CHECK: 手写中“…”与“2n-1”之间有一处被涂黑划去的字迹无法辨认；前几个冒号形似等号
通过连续相等位移所用时间之比
\[ t_1:t_2:t_3:\cdots:t_n=1:\sqrt{2}-1:\sqrt{3}-\sqrt{2}:\sqrt{n+1}-\sqrt{n} \] % CHECK: 末项手写为 sqrt(n+1)-sqrt(n)（通常结论的末项为 sqrt(n)-sqrt(n-1)），照原样转录
%%%END
%%%PAGEEND 66
